% !TeX encoding = UTF-8
% Español de España. Cargar un único arquetipo desde el archivo principal.
\newcommand{\cvsummary}{Ingeniero mecánico por NYU Abu Dhabi, actualmente cursando un máster en Mecatrónica, con experiencia práctica en sistemas eléctricos industriales, automatización y digitalización. Experiencia en diagnóstico y resolución de averías, mejora de la eficiencia operativa y desarrollo de herramientas digitales para procesos de ingeniería.}

% --- COMPETENCIAS ---
\newcommand{\skills}{
    \section*{COMPETENCIAS TÉCNICAS}
    {\small
    \textbf{Control y automatización:} Control PID con anti-windup y gestión de saturación de actuadores, fundamentos de espacio de estados y LQR, máquinas de estados finitos, programación de PLC, control de motores de corriente continua, control embebido en tiempo real, MATLAB/Simulink.

    \vspace{1pt}
    \textbf{Sistemas embebidos y electrónica:} Microcontroladores (Arduino/C++), FPGA, PWM, ADC/DAC, codificadores en cuadratura, controladores de motores, IMU y fusión sensorial, Raspberry Pi, IoT, sensores e instrumentación, diseño lógico digital (Quartus), integración de hardware y software.

    \vspace{1pt}
    \textbf{Sistemas eléctricos:} Esquemas unifilares, análisis de subestaciones de alta tensión (armónicos), corrección del factor de potencia (baterías de condensadores), AutoCAD Electrical, sistemas industriales de gestión de almacenes (SGA) y trazabilidad (Industria 4.0).

    \vspace{1pt}
    \textbf{Diseño mecánico y simulación:} SolidWorks (CSWA), Solid Edge, CATIA V5, Siemens NX, CREO, ANSYS (Mechanical, Fluent), análisis por elementos finitos (FEA), dinámica de fluidos computacional (CFD), diseño de sistemas térmicos (TEC y climatización).

    \vspace{1pt}
    \textbf{Fabricación y calidad:} CNC, fabricación aditiva y FDM, inyección de plásticos (análisis con Moldflow), materiales compuestos, tolerancias geométricas y dimensionales (GD\&T), control de calidad ISO 9001, Lean Six Sigma White Belt.

    \vspace{1pt}
    \textbf{Programación:} C++, Python, MATLAB, comunicación Qt/TCP con controladores de motores, JavaScript, SQL, Git.

    \vspace{1pt}
    \textbf{Gestión de proyectos:} MS Project, Jira, Notion, coordinación multidisciplinar e informes técnicos.

    \vspace{1pt}
    \textbf{Idiomas:} Español nativo, inglés C2, francés C2, italiano A1 y alemán A1.
    }
}   

% --- PROYECTOS ---
\newcommand{\vawttitle}{Aerogenerador de eje vertical: aprovechamiento de energía eólica urbana}
\newcommand{\vawtdates}{09/2024 - 05/2025}
\newcommand{\vawtlocation}{Abu Dabi, EAU}
\newcommand{\vawtdescOne}{Desarrollé un nuevo diseño de aerogenerador de eje vertical con palas de arrastre y sustentación para el viento urbano de EAU, aplicando diseño mecánico y simulaciones CFD para optimizar el aprovechamiento energético. Patente en trámite.}
\newcommand{\vawtdescTwo}{Evalué la aerodinámica del rotor mediante el modelado del sistema de aprovechamiento eólico y el análisis de los resultados de simulaciones CFD en ANSYS Fluent.}

\newcommand{\fsaetitle}{Formula SAE -- Suspensión y chasis de fibra de carbono}
\newcommand{\fsaedates}{Otoño de 2024 -- Actualidad}
\newcommand{\fsaelocation}{Abu Dabi, EAU}
\newcommand{\fsaedescOne}{Desarrollé sistemas estructurales del vehículo, incluidos la suspensión y el chasis de fibra de carbono, combinando modelado en SolidWorks y análisis por elementos finitos en ANSYS.}
\newcommand{\fsaedescTwo}{Evalué el movimiento de la suspensión mediante simulaciones y análisis cinemáticos del mecanismo en SolidWorks y ANSYS.}

\newcommand{\tactitle}{Mininevera -- Sistema termoeléctrico TEC de dos etapas}
\newcommand{\tacdates}{Primavera de 2025}
\newcommand{\taclocation}{Abu Dabi, EAU}
\newcommand{\tacdescOne}{Desarrollé un sistema de refrigeración termoeléctrica de dos etapas sin refrigerante, sustituyendo el circuito frigorífico por módulos TEC.}
\newcommand{\tacdescTwo}{Integré sensores e instrumentación en tiempo real en el conjunto de refrigeración para supervisar su comportamiento térmico durante el funcionamiento.}

\newcommand{\orbitaltitle}{OrbitalGrid -- Sistema de información geográfica espacial}
\newcommand{\orbitaldates}{06/2024 - Actualidad}
\newcommand{\orbitallocation}{Abu Dabi, EAU}
\newcommand{\orbitaldescOne}{Desarrollé un sistema de integración de datos espaciales mediante un proyecto de cartografía SIG basado en SQL que combina escaneado 3D y datos satelitales.}
\newcommand{\orbitaldescTwo}{Implementé la gestión de datos geoespaciales en PostgreSQL, estructurando la información para aplicaciones de cartografía aeroespacial.}

\newcommand{\softrobtitle}{Exoesqueleto robótico para rehabilitación tras un ictus}
\newcommand{\softrobdates}{04/2024 - 05/2024}
\newcommand{\softroblocation}{Abu Dabi, EAU}
\newcommand{\softrobdescOne}{Diseñé un mecanismo de rehabilitación accionado mediante actuadores, basado en un exoesqueleto de dedos inspirado en el origami que integra hilos de aleación con memoria de forma en una estructura flexible.}
\newcommand{\softrobdescTwo}{Integré materiales blandos y elementos activos mediante un diseño de polímero Dragon Skin e hilos SMA, combinando interfaces flexibles y accionamiento por respuesta térmica.}

\newcommand{\exoplanettitle}{Simulador de tránsitos de exoplanetas}
\newcommand{\exoplanetdates}{Primavera de 2025}
\newcommand{\exoplanetlocation}{Abu Dabi, EAU}
\newcommand{\exoplanetdescOne}{Desarrollé un modelo físico parametrizado de tránsitos de exoplanetas en Python mediante la implementación de mecánica orbital y física estelar.}
\newcommand{\exoplanetdescTwo}{Generé datos de observación sintéticos en forma de series temporales de curvas de luz, calculando el comportamiento de los tránsitos a partir de parámetros astrofísicos.}

\newcommand{\windTurbinetitle}{Aerogenerador marino flotante -- Estabilización activa}
\newcommand{\windTurbinedates}{05/2026 - 06/2026}
\newcommand{\windTurbinelocation}{Besançon, Francia}
\newcommand{\windTurbinedescOne}{Diseñé y construí en SolidWorks un prototipo de aerogenerador semisumergible de tres columnas a escala 1:1000, estableciendo su estabilidad en cabeceo y balanceo con GM = 5,38 cm; integré un MPU6050 y una suspensión cardán accionada por servomotores para estabilizar activamente la góndola en ensayos en canal de oleaje.}
\newcommand{\windTurbinedescTwo}{Establecí la estabilidad pasiva en cabeceo, con GM = 5,38 cm y rigidez de restitución K = 0,0607 N$\cdot$m/rad, mediante el cálculo hidrostático de un prototipo semisumergible de tres columnas a escala 1:1000.}

\newcommand{\carbonmetertitle}{CARBONMETRE -- Ecodiseño y medición de carbono}
\newcommand{\carbonmeterdates}{09/2025 - 12/2025}
\newcommand{\carbonmeterlocation}{Besançon, Francia}
\newcommand{\carbonmeterdescOne}{Apliqué principios metrológicos a la medición de carbono en el proyecto de ingeniería CARBONMETRE, combinando medición, fabricación y ecodiseño.}
\newcommand{\carbonmeterdescTwo}{Relacioné la fabricación con la ingeniería ambiental mediante la aplicación de principios de fabricación industrial a un proyecto de medición de carbono.}

\newcommand{\hvactitle}{Acondicionamiento térmico de una habitación de residencia -- Refrigeración pasiva y evaporativa}
\newcommand{\hvacdates}{06/2026}
\newcommand{\hvaclocation}{Besançon, Francia}
\newcommand{\hvacdescOne}{Reduje hasta 3 \textdegree{}C la temperatura de una habitación de 9 m$^2$ respecto a las condiciones previas a la intervención, mediante una pantalla solar exterior y un refrigerador de agua congelada con ventilador.}
\newcommand{\hvacdescTwo}{Bloqueé la exposición solar en las horas punta de 14:00 a 21:00 mediante capas exteriores de papel de aluminio reflectante, EPS de 1 cm y cartón ondulado.}

\newcommand{\controlBoxtitle}{Instalación eléctrica integral -- PLC para eventos de mediana escala}
\newcommand{\controlBoxdates}{06/2023 - 08/2023}
\newcommand{\controlBoxlocation}{Ciudad de México, México}
\newcommand{\controlBoxdescOne}{Reduje un 30\% el tiempo de montaje del sistema eléctrico respecto a la instalación anterior mediante el diseño de un cuadro de control con PLC de 12 entradas para eventos.}
\newcommand{\controlBoxdescTwo}{Desarrollé un sistema de control eléctrico basado en PLC para instalaciones de eventos de mediana escala, adaptado a los requisitos de funcionamiento de los equipos audiovisuales.}

\newcommand{\hyperlooptitle}{NYC Hyperloop -- Vehículo a escala reducida}
\newcommand{\hyperloopdates}{01/2024 - 05/2024}
\newcommand{\hyperlooplocation}{Nueva York, EE. UU.}
\newcommand{\hyperloopdescOne}{Lideré el desarrollo del sistema mecánico como responsable del subequipo mecánico de un proyecto Hyperloop a escala reducida para un concepto de trayecto Nueva York--Long Island.}
\newcommand{\hyperloopdescTwo}{Evalué el comportamiento de la levitación magnética mediante el modelado multifísico del sistema de transporte y simulaciones en COMSOL y ANSYS.}

\newcommand{\dronetitle}{Cuadricóptero funcional -- Estructura ligera}
\newcommand{\dronedates}{03/2024 - 05/2024}
\newcommand{\dronelocation}{Nueva York, EE. UU.}
\newcommand{\dronedescOne}{Reduje un 40\% el peso del cuadricóptero respecto al diseño de referencia del proyecto mediante una estructura fabricada por impresión 3D en PLA y TPU.}
\newcommand{\dronedescTwo}{Evalué la resistencia al impacto de la estructura del cuadricóptero mediante ensayos de caída contrastados con los criterios de rendimiento durante el desarrollo.}

\newcommand{\trentxwbtitle}{Rolls-Royce Trent XWB -- Modelo 3D funcional}
\newcommand{\trentxwbdates}{01/2024 - 05/2024}
\newcommand{\trentxwblocation}{Nueva York, EE. UU.}
\newcommand{\trentxwbdescOne}{Codesarrollé un conjunto CAD integrado de propulsión mediante el modelado de los componentes del motor en SolidWorks y Fusion 360, obteniendo un modelo 3D funcional completo del Trent XWB.}
\newcommand{\trentxwbdescTwo}{Evalué el comportamiento aerodinámico del modelo de motor mediante análisis y simulaciones en ANSYS.}

\newcommand{\liostitle}{LIOS -- Asistente de voz con IA}
\newcommand{\liosdates}{12/2024 - Actualidad}
\newcommand{\lioslocation}{Abu Dabi, EAU}
\newcommand{\liosdescOne}{Desarrollé un asistente de voz con IA en Python para automatizar tareas mediante la integración de procesamiento del lenguaje natural, datos en tiempo real y API externas.}
\newcommand{\liosdescTwo}{Habilité interacciones sensibles al contexto mediante memoria contextual que conserva la información relevante en los procesos del asistente.}

\newcommand{\fitpaltitle}{FitPal -- Dispositivo inteligente para entrenamiento físico}
\newcommand{\fitpaldates}{09/2021 - 12/2021}
\newcommand{\fitpallocation}{Abu Dabi, EAU}
\newcommand{\fitpaldescOne}{Codesarrollé un dispositivo embebido de medición basado en M5Stack para analizar ejercicios en tiempo real, integrando sensores y procesamiento en un producto de entrenamiento físico.}
\newcommand{\fitpaldescTwo}{Conecté el dispositivo a una interfaz móvil mediante la integración de Arduino IoT App y funciones de transferencia de datos IoT.}

\newcommand{\covidtitle}{Cabina de desinfección automatizada}
\newcommand{\coviddates}{09/2020 - 02/2021}
\newcommand{\covidlocation}{Ciudad de México, México}
\newcommand{\coviddescOne}{Desarrollé un prototipo de cabina de desinfección controlado por Arduino, integrando la electrónica de control y un sistema de nebulización.}
\newcommand{\coviddescTwo}{Añadí comunicación inalámbrica al sistema mediante la conexión sin cables de los elementos de control del prototipo.}

\newcommand{\rovertitle}{Rover teleoperado -- Escaneado espacial 3D}
\newcommand{\roverdates}{09/2026 - Actualidad}
\newcommand{\roverlocation}{Oviedo, España}
\newcommand{\roverdescOne}{Desarrollo una plataforma móvil de sensores teleoperada, combinando control embebido, adquisición 3D y visualización remota como objetivo de integración en un rover.}
\newcommand{\roverdescTwo}{Desarrollo el control embebido del rover mediante la integración del subsistema de control, con el funcionamiento por órdenes remotas como hito objetivo.}

\newcommand{\microsensortitle}{Microsensor piezorresistivo de presión de detonación}
\newcommand{\microsensordates}{05/2026 - 06/2026}
\newcommand{\microsensorlocation}{Besançon, Francia}
\newcommand{\microsensordescOne}{Diseñé un microsensor de presión piezorresistivo, elaboré instrucciones detalladas de fabricación que abarcan pulverización catódica y lift-off, y automaticé en Python los cálculos geométricos con restricciones de factor de galga, $\Delta R/R_0$ y resistividad.}
\newcommand{\microsensordescTwo}{Lideré la selección de materiales bajo restricciones eléctricas y mecánicas para un diseño con interfaz de 50 $\Omega$ y resistencia a presiones de 10 GPa, equilibrando requisitos eléctricos, mecánicos y económicos a escala de oblea.}

\newcommand{\missiletitle}{Simulador de dinámica de interceptación de misiles}
\newcommand{\missilelocation}{Besançon, Francia}
\newcommand{\missiledescOne}{Desarrollé en MATLAB una simulación 3D de interceptación de misiles mediante guiado por navegación proporcional, integración numérica y visualización de trayectorias en tiempo real con detección de impacto.}

% --- EXPERIENCIA ---
\newcommand{\carbonorg}{Energy Innovation Capital}
\newcommand{\carbontitle}{Analista de mercado de captura de carbono}
\newcommand{\carbondates}{10/2023 - 01/2024}
\newcommand{\carbonlocation}{Abu Dabi, EAU}
\newcommand{\carbonbullets}{%
  \item Apoyé decisiones de inversión mediante recomendaciones estratégicas sobre empresas emergentes CCUS, analizando tecnologías de captura de carbono para un fondo de capital riesgo de Abu Dabi.
  \item Traduje investigación técnica para responsables de la toma de decisiones mediante presentaciones a analistas sénior que sintetizaban avances científicos en captura de carbono.
  \item Identifiqué los factores impulsores de la tecnología y el mercado mediante análisis de tendencias que combinaban investigación técnica y datos sectoriales.
  \item Contribuí a la evaluación del riesgo de inversión del fondo mediante el análisis de los avances tecnológicos CCUS junto con las condiciones de mercado.
}

\newcommand{\spacelaborg}{NYU Abu Dhabi, Centro de Ciencias Espaciales y Astrofísica}
\newcommand{\spacelabtitle}{Ingeniero e investigador de I+D -- Plataforma de ciencias planetarias}
\newcommand{\spacelabdates}{05/2025 - 09/2025}
\newcommand{\spacelablocation}{Abu Dabi, EAU}
\newcommand{\spacelabbullets}{%
  \item Mejoré un 40\% la eficiencia del procesamiento de imágenes respecto al proceso anterior mediante la optimización de la compresión de datos, la reproyección cartográfica y la corrección de imágenes asistida por aprendizaje automático.
  \item Transformé imágenes de una sonda espacial en una \href{https://romvio.github.io/mars-viewer/}{aplicación web} interactiva con React y JavaScript, facilitando el acceso a datos de la misión HOPE mediante una interfaz cartográfica visual; adapté la plataforma a iOS y la presenté a la App Store de Apple en EAU.
  \item Integré imágenes de distintas pasadas orbitales en mapas planetarios coherentes mediante alineación de solapamientos, corrección de diferencias de brillo e implementación de proyecciones cilíndrica y Robinson.
}

\newcommand{\geoorg}{Global Education NYUAD}
\newcommand{\geotitle}{Coordinador de Global Education -- Estudiantes visitantes}
\newcommand{\geodates}{12/2025 - 05/2025}
\newcommand{\geolocation}{Abu Dabi, EAU}
\newcommand{\geobullets}{%
  \item Coordiné las operaciones de programas sobre el terreno en EAU mediante la organización de salidas académicas y culturales para estudiantes visitantes de NYU Abu Dhabi.
  \item Organicé la ejecución de las actividades mediante el liderazgo del equipo de coordinadores y la planificación logística conjunta.
  \item Contribuí a la seguridad de las actividades en grupo mediante protocolos y medidas preventivas incorporados a la logística de las salidas.
  \item Mantuve el apoyo a los participantes como enlace designado, coordinando la comunicación entre los estudiantes y los servicios universitarios.
}

\newcommand{\geoOneorg}{Global Education NYUAD}
\newcommand{\geoOnetitle}{Coordinador de Global Education -- Estudiantes de primer curso}
\newcommand{\geoOnedates}{09/2026 - 12/2024}
\newcommand{\geoOnelocation}{Abu Dabi, EAU}
\newcommand{\geoOnebullets}{%
  \item Facilité la incorporación de estudiantes de primer curso mediante sesiones de orientación e introducción a la vida universitaria.
  \item Ofrecí orientación continuada mediante mentoría y asesoramiento entre estudiantes para facilitar el uso de los recursos del campus.
  \item Comuniqué información práctica sobre las instalaciones y los servicios universitarios mediante visitas al campus.
  \item Contribuí a la integración de los grupos mediante la planificación de eventos académicos y culturales de inclusión para estudiantes de primer curso.
}

\newcommand{\hcorg}{Centro de Salud de NYUAD}
\newcommand{\hctitle}{Empleo de estudiante -- Responsable de redes sociales}
\newcommand{\hcdates}{09/2024 - 05/2025}
\newcommand{\hclocation}{Abu Dabi, EAU}
\newcommand{\hcbullets}{%
  \item Promoví la participación del alumnado mediante campañas de bienestar y comunicación sanitaria adaptadas a la comunidad de NYU Abu Dhabi.
  \item Elaboré materiales de comunicación digital y contenidos para redes sociales destinados a promover iniciativas de salud.
  \item Actué como contacto principal para consultas del alumnado, coordinando las respuestas a preguntas dirigidas al centro de salud.
  \item Contribuí a la resolución de incidencias mediante la gestión de consultas, combinando atención directa al alumnado y comunicación del centro de salud.
}

\newcommand{\alpesorg}{Alpes Technologies México}
\newcommand{\alpestitle}{Ingeniero de proyectos júnior -- SGA y sistemas eléctricos}
\newcommand{\alpesdates}{06/2024 - 09/2024}
\newcommand{\alpeslocation}{Ciudad de México, México}
\newcommand{\alpesbullets}{%
  \item Fabriqué e instalé baterías de condensadores industriales con filtros de armónicos mediante el análisis de armónicos y picos de tensión en subestaciones, alcanzando factores de potencia de hasta 0,99 y reduciendo los costes energéticos un 15\% y las paradas eléctricas un 25\%.
  \item Sustituí el seguimiento manual de inventario por un sistema de gestión de almacenes con códigos QR desarrollado en HTML/JavaScript y Google Drive, mejorando un 60\% la eficiencia operativa y de fabricación.
  \item Analicé daños eléctricos tras el incendio de un complejo hotelero mediante mediciones in situ y presenté recomendaciones correctivas a la alta dirección.
}

\newcommand{\epigsaorg}{EPIGSA}
\newcommand{\epigsatitle}{Asistente de ingeniería de datos -- SIG e hidráulica}
\newcommand{\epigsadates}{05/2024 - 07/2024}
\newcommand{\epigsalocation}{Ciudad de México, México}
\newcommand{\epigsabullets}{%
  \item Reduje un 25\% los errores en datos SIG mediante la automatización de controles de calidad para proyectos municipales de infraestructuras de tratamiento de aguas.
  \item Mantuve y validé registros de infraestructuras de ingeniería en PostgreSQL, garantizando la coherencia y accesibilidad de los datos para equipos multidisciplinares.
}

\newcommand{\solarshiporg}{Solar Ship}
\newcommand{\solarshiptitle}{Ingeniero en prácticas de estrategia y diseño de movilidad aérea}
\newcommand{\solarshipdates}{05/2023 - 08/2023}
\newcommand{\solarshiplocation}{Toronto, Canadá}
\newcommand{\solarshipbullets}{%
  \item Desarrollé conceptos de logística aérea para zonas remotas de Zambia y Tanzania mediante la planificación de infraestructuras para comunidades sin conexión a la red.
  \item Contribuí al diseño de un sistema aeronáutico totalmente eléctrico para la iniciativa Peace+Freedom, orientado a mejorar la conectividad de las comunidades.
  \item Relacioné el diseño del vehículo con los requisitos de infraestructura mediante un concepto de logística aérea que consideraba las necesidades de acceso junto con los sistemas de la aeronave.
  \item Integré objetivos técnicos y sociales en un proyecto multinacional de movilidad, combinando ingeniería aeroespacial y requisitos de acceso de las comunidades.
}

\newcommand{\crousorg}{CROUS Besançon Bourgogne Franche-Comté}
\newcommand{\croustitle}{Empleo de estudiante -- Restauración universitaria}
\newcommand{\crousdates}{10/2024 - Actualidad}
\newcommand{\crouslocation}{Besançon, Francia}
\newcommand{\crousbullets}{%
  \item Mantuve la continuidad del trabajo de cocina en horas punta mediante apoyo a la preparación de comidas universitarias de gran volumen.
  \item Apliqué procedimientos de limpieza de cocina y asumí tareas de higiene y seguridad alimentaria.
  \item Contribuí a la continuidad del servicio mediante la reposición de suministros y el apoyo durante la actividad.
  \item Colaboré en tareas coordinadas de cocina y servicio, adaptándome a necesidades operativas con plazos ajustados.
}

\newcommand{\rallyorg}{Campeonato Mundial de Rally -- Guanajuato}
\newcommand{\rallytitle}{Comisario técnico oficial de la FIA}
\newcommand{\rallydates}{11/2019 - 04/2020}
\newcommand{\rallylocation}{Guanajuato, México}
\newcommand{\rallybullets}{%
  \item Realicé inspecciones técnicas reglamentarias y el marcaje oficial de neumáticos en una prueba del Campeonato Mundial de Rally de la FIA, siguiendo los procedimientos de verificación establecidos.
  \item Contribuí al cumplimiento técnico de los vehículos de competición mediante inspecciones conforme a los requisitos de la prueba.
  \item Coordiné inspecciones con comisarios responsables y pilotos bajo la presión operativa de una competición en curso.
  \item Apliqué procedimientos establecidos de marcaje y verificación en trabajos de inspección descritos conforme a normas ISO 9001.
}

\newcommand{\papaloteorg}{Papalote Museo del Niño}
\newcommand{\papalotetitle}{Auxiliar de atención al público}
\newcommand{\papalotedates}{05/2019 - 07/2019}
\newcommand{\papalotelocation}{Ciudad de México, México}
\newcommand{\papalotebullets}{%
  \item Expliqué sistemas técnicos a personas no especializadas mediante visitas a exposiciones de ciencia y tecnología, adaptando las explicaciones a niños y adultos.
  \item Contribuí a mejoras de servicio implantadas mediante propuestas prácticas para la experiencia de los visitantes y el funcionamiento del museo.
  \item Atendí a visitantes de distintas edades, adaptando la comunicación a sus necesidades.
  \item Relacioné la ciencia con la participación práctica mediante demostraciones y explicaciones accesibles de conceptos técnicos en las exposiciones.
}
